% CPSY 1291 -- Recitation handout 11
% Compile:  latexmk -lualatex handout-11-attention-transformers-generative.tex
\documentclass[11pt]{article}
\usepackage{cpsy1291-handout}

\begin{document}

\handouthead{11}{Attention and the transformer}
  {What Assignment 4 needs from Lecture 17}
  {Week 11 \ $\cdot$\ hold this after Lecture 17; Assignment 4 is due Tue 11/24}

\section*{What this session is for}
\addcontentsline{toc}{section}{What this session is for}

This is the last handout about content. Assignment 4 covers Lectures 15--17 and
is due on Tuesday 11/24. Recitation 10 did recurrence and dynamics (Lectures 15
and 16); this handout finishes the assignment's scope with Lecture 17. Generative
models (Lecture 14) and large language models (Lecture 18) are not on the
assignment; they are on Exam 2, and Recitation 12 reviews them.

Read it in two passes. Section~1 (attention) is the part everything else depends
on, and it is worth doing slowly with a pencil. Sections~2 and 3 are short.

\begin{keybox}
\ask{The one idea.} Attention computes a set of \emph{weights} that say how much
each position should contribute, and takes a weighted average of what the
positions offer. A modern Hopfield update does exactly the same thing, with the
cue as the query and the stored patterns as keys and values (Section~3). The
differences are in what is averaged and what determines the weights.
\end{keybox}

\section{Attention}

\subsection{Tokens are a matrix}

A sequence of $T$ tokens, each embedded as a $d$-dimensional vector, is a matrix
$\boldsymbol{X}$ of shape $(T, d)$. Everything in a transformer operates on that
matrix. Note what has been given up compared to an RNN: there is no order in a
matrix's rows unless you put it there, which is why a learned table of
\textbf{positional embeddings} is added to each symbol's vector. What is gained:
any position reaches any other in one step, and all $T$ rows come out of a few
matrix products. The price is $T\times T$ scores, so doubling the sequence
length quadruples the cost.

\subsection{Q, K, V are three views of the same input}

\begin{lstlisting}
Q = X @ Wq        # (T, d_k)   what each position is looking for
K = X @ Wk        # (T, d_k)   what each position offers
V = X @ Wv        # (T, d_v)   what each position would contribute
\end{lstlisting}

All three are linear projections of the \emph{same} $\boldsymbol{X}$ --- that is
what makes it \emph{self}-attention. The names are a retrieval metaphor: a query
is matched against keys, and the matching keys' values are returned.

\subsection{Scaled dot-product attention}

\[
\mathrm{Attention}(\boldsymbol{Q},\boldsymbol{K},\boldsymbol{V})
= \mathrm{softmax}\!\left(\frac{\boldsymbol{Q}\boldsymbol{K}^{\!\top}}{\sqrt{d_k}}\right)\boldsymbol{V}
\]

Masked positions get a score of $-\infty$ before the softmax, so their weight is exactly 0. Writing this as code is part of Assignment~4.

Track the shapes, because they are the whole thing:

\begin{center}
\begin{tabular}{@{}lll@{}}
\toprule
Step & Shape & Meaning \\
\midrule
$\boldsymbol{Q}\boldsymbol{K}^{\!\top}$ & $(T, T)$ & every position against every position \\
after softmax & $(T, T)$ & each \emph{row} is a distribution over positions \\
$\boldsymbol{A}\boldsymbol{V}$ & $(T, d_v)$ & each position: a weighted average of all values \\
\bottomrule
\end{tabular}
\end{center}

\begin{keybox}
\ask{Three details that are not decoration.}
\textbf{The $\sqrt{d_k}$}: dot products of $d_k$-dimensional vectors grow like
$\sqrt{d_k}$, and a softmax of large numbers saturates to a one-hot vector with
almost no gradient. Remove it and training fails for large $d_k$.
\textbf{\texttt{dim=-1}}: the softmax runs across \emph{keys}, so each row is a
distribution over what to attend to. Softmax over the wrong axis produces
nonsense that runs perfectly.
\textbf{The mask is $-\infty$, not $0$}: it must be applied \emph{before} the
softmax, because zeroing afterwards leaves rows that no longer sum to one.
\end{keybox}

\subsection{Multi-head: split, attend, concatenate}

Split the $d$ dimensions into $H$ groups, run the same operation on each,
concatenate the results, and let a final matrix $\boldsymbol{W}_O$ mix the heads
back to width $d$. Each head has its own three projections, so different heads
can attend to different things at the cost of one wide head.

\begin{lstlisting}
# (N, T, d) -> (N, H, T, d//H)
Q = Q.view(N, T, H, d // H).transpose(1, 2)
# ... attention exactly as above, batched over the H axis ...
out = out.transpose(1, 2).reshape(N, T, d)     # concatenate the heads back
\end{lstlisting}

The \texttt{transpose(1, 2)} before and after is the only subtle part. It puts
the head axis next to the batch axis so that the matrix multiplies treat heads
as independent problems.

\subsection{Reading an attention map, carefully}

The $(T, T)$ matrix $\boldsymbol{A}$ is tempting to plot and interpret. Single
heads in trained language models do track relations people can name --- verb to
object, pronoun to referent --- but three cautions apply:

\begin{itemize}
  \item There are $H$ heads per layer and many layers, so ``the attention map''
        is one choice among many.
  \item A map shows \emph{where} a head reads, not \emph{why} the output is what
        it is.
  \item A map is a hypothesis. Changing the input, or switching the head off,
        tests it.
\end{itemize}

\begin{keybox}
\ask{The one-sentence version.} An attention map is a hypothesis about where
information is read, and needs an intervention before it explains anything ---
the discipline of the attribution maps in Recitation 9.
\end{keybox}

\section{Training and generation}

\subsection{One forward pass, $T$ training examples}

A language model predicts the next symbol at every position. The target at
position $t$ is the symbol at $t+1$, and the loss is the cross-entropy averaged
over positions, in nats:
\[
\mathcal{L} = -\frac{1}{T}\sum_{t=1}^{T}\log \hat{y}^{\langle t\rangle}_{\,y^{\langle t+1\rangle}}.
\]
The \textbf{causal mask} sets every score above the diagonal to $-\infty$, so
position $t$ cannot read the symbol it must predict. With the mask, one forward
pass gives all $T$ predictions. The input is the true text, not the model's own
guesses: this is teacher forcing, from Recitation 10.

\subsection{The bigram baseline}

Count, in the training text, how often each symbol follows each other one:
\[
P(\text{next}=b \mid \text{current}=a) = \frac{\text{count}(a\to b)}{\text{count}(a)}.
\]
Score the table with the same cross-entropy, on the same held-out text. A bigram
table reads one symbol of context and the transformer can read the whole left
context, so the gap between the two losses is what the longer context is worth.
As in the baselines lecture, a score means something only next to a baseline
computed on the same data.

\subsection{Generation is a loop}

\begin{lstlisting}
for _ in range(n_new):
    logits = model(idx)[:, -1, :] / temperature      # last position only
    probs = torch.softmax(logits, dim=-1)
    nxt   = torch.multinomial(probs, num_samples=1)  # SAMPLE, do not argmax
    idx   = torch.cat([idx, nxt], dim=1)
\end{lstlisting}

If the text grows longer than the context, keep only the last $T$ symbols.

\begin{center}
\begin{tabular}{@{}ll@{}}
\toprule
\texttt{temperature} & What it does \\
\midrule
$\to 0$ & approaches greedy: repetitive, often loops \\
$= 1$ & samples from the model's own distribution \\
$> 1$ & flatter, more surprising, less coherent \\
\bottomrule
\end{tabular}
\end{center}

\begin{keybox}
\ask{Sampling is not a detail of the model.} Two very different outputs from the
same weights can differ only in the temperature, so any claim about ``what the
model does'' has to state it.
\end{keybox}

\section{The Hopfield update is one attention step}

A modern Hopfield network keeps its stored patterns as the rows of
$\boldsymbol{X}$ and recalls from a cue $\boldsymbol{s}$ with a softmax:
\[
\boldsymbol{s}\;\leftarrow\;\boldsymbol{X}^{\top}\,\mathrm{softmax}\big(\beta\,\boldsymbol{X}\,\boldsymbol{s}\big).
\]
Read it in three steps: $\boldsymbol{X}\boldsymbol{s}$ gives one match score per
stored pattern, the softmax turns the scores into weights, and
$\boldsymbol{X}^{\top}$ times the weights is a weighted average of the stored
patterns. That is attention for one query,
$\mathrm{softmax}(\boldsymbol{q}\boldsymbol{K}^{\!\top}/\sqrt{d_k})\boldsymbol{V}$,
with the cue as the query, the stored patterns as both keys and values (no
learned projections: $\boldsymbol{W}_K=\boldsymbol{W}_V=\boldsymbol{I}$), and
$\beta$ in the role of $1/\sqrt{d_k}$. A large $\beta$ sharpens the weights, so
recall picks out one stored pattern.

\section{Exercises}

\ask{1.} You delete the $\sqrt{d_k}$ from scaled dot-product attention. Training
works at $d_k = 8$ and fails at $d_k = 512$. Why?

\ask{2.} You apply a causal mask by multiplying the attention weights by 0
\emph{after} the softmax. What is wrong?

\ask{3.} You take softmax over \texttt{dim=-2} instead of \texttt{dim=-1}. Does
anything error? What does the model now compute?

\ask{4.} A colleague shows an attention map and says it proves the model bases
its answer on the word ``not''. Give two objections, and the test that would
settle it.

\ask{5.} Generation at \texttt{temperature=0.1} produces the same sentence
repeated forever. Explain in terms of the sampling loop.

\ask{6.} In the toy text ``that hat ate the tea'', what is
$P(\text{next}=\text{a}\mid\text{current}=\text{h})$ under a bigram table? Your
transformer's held-out loss equals the table's. What has it learned?

\ask{7.} In the modern Hopfield update, which quantity plays the query, which the
keys, which the values, and which $1/\sqrt{d_k}$?

\subsection*{Answers}

\ask{1.} Dot products of $d_k$-dimensional vectors have magnitude growing like
$\sqrt{d_k}$. At $d_k = 512$ the unscaled scores are large, the softmax saturates
to nearly one-hot, and its gradient is nearly zero, so nothing trains. At
$d_k = 8$ the scores are small enough that it hardly matters.

\ask{2.} The softmax has already distributed weight to the masked positions, so
zeroing afterwards leaves rows that do not sum to one --- and the unmasked
positions get less than their share. Mask with $-\infty$ \emph{before} the
softmax.

\ask{3.} Nothing errors: the matrix is $(T, T)$, so both axes are valid. You now
normalize down the columns, so each row is no longer a distribution over
positions and the output is a weighted combination with weights that do not sum
to one. It trains, badly, with no diagnostic.

\ask{4.} (i) There are many heads and many layers; ``the attention map'' is one
choice among many. (ii) A map shows where a head reads, not why the output is
what it is. The test is an intervention: change the input (remove or replace
``not'') or switch the head off, and see whether the answer changes.

\ask{5.} At low temperature the softmax is nearly one-hot, so sampling is
effectively \texttt{argmax}. The model picks its most likely continuation every
time, and always taking the most likely symbol tends to loop. Raising the
temperature breaks the loop.

\ask{6.} The symbol h is followed by a twice and by e once, so $2/3$. A
transformer that only matches the bigram table has used no more than one symbol
of context: its attention is not reading the rest of the left context. Check the
mask, the positional embeddings, and whether it trained at all.

\ask{7.} The cue $\boldsymbol{s}$ is the query; the stored patterns
$\boldsymbol{X}$ are both the keys and the values; $\beta$ plays $1/\sqrt{d_k}$.

\section*{Cheat sheet}
\addcontentsline{toc}{section}{Cheat sheet}

\begin{center}
\begin{tabular}{@{}p{0.44\textwidth}p{0.5\textwidth}@{}}
\toprule
\textbf{You want} & \textbf{You write} \\
\midrule
attention scores          & \texttt{Q @ K.transpose(-2,-1) / sqrt(d\_k)} \\
a causal mask             & \texttt{masked\_fill(mask == 0, float('-inf'))} \emph{before} softmax \\
the softmax axis          & \texttt{dim=-1} --- over keys, always \\
split into heads          & \texttt{view(N,T,H,d//H).transpose(1,2)} \\
the next-symbol target    & the input shifted by one position \\
a bigram baseline         & \texttt{count(a$\to$b) / count(a)}, scored on held-out text \\
sample a token            & \texttt{torch.multinomial(probs, 1)} --- not \texttt{argmax} \\
a modern Hopfield update  & \texttt{X.T @ softmax(beta * X @ s)} \\
\bottomrule
\end{tabular}
\end{center}

\end{document}
